\chapter{Phương pháp đề xuất}
\chaptergoal{Xây dựng pipeline recommendation hoàn chỉnh từ dữ liệu chuẩn hóa đến artifact phục vụ suy luận, đồng thời đảm bảo khả năng tái lập thực nghiệm.}

\section{Kiến trúc phương pháp}
Thay vì đặt toàn bộ logic trong một module duy nhất, đề tài sử dụng kiến trúc ba tầng tách biệt:
\begin{itemize}
    \item \textbf{ETL layer}: thu thập và chuẩn hóa dữ liệu vào \texttt{core\_*};
    \item \textbf{Training layer}: sinh profile, huấn luyện mô hình, đánh giá và xuất artifacts;
    \item \textbf{Serving layer}: nạp artifacts để cung cấp API recommendation.
\end{itemize}

Cấu trúc này giải quyết hai vấn đề thực tế: (i) tránh trộn lẫn code offline và code runtime; (ii) cho phép thực hiện nhiều vòng lặp thí nghiệm mà không ảnh hưởng dịch vụ discovery đang chạy.

\section{Định nghĩa đơn vị biểu diễn: movie profile}
Mỗi phim được ánh xạ thành một chuỗi văn bản tổng hợp gọi là \textit{movie profile}. Trong phiên bản hiện tại:
\[
\text{profile} = \text{title} \oplus \text{overview} \oplus \text{genres} \oplus \text{review snippets}
\]
trong đó $\oplus$ là phép nối có cấu trúc. Lý do lựa chọn thiết kế này là cân bằng giữa \textit{độ giàu thông tin} và \textit{chi phí tính toán}. Nếu chỉ dùng \textit{overview}, mô hình mất đi tín hiệu cảm nhận người dùng; nếu dùng toàn bộ review thô, hồ sơ phim trở nên quá dài và nhiễu.

\section{Pipeline thuật toán}
\begin{enumerate}[label=\arabic*.]
    \item Nạp dữ liệu từ \texttt{core\_movies}, \texttt{core\_reviews}, \texttt{core\_genres}.
    \item Lọc và chuẩn hóa văn bản review theo chính sách ngôn ngữ.
    \item Sinh \textit{movie profile} cho từng phim.
    \item Vector hóa profile theo từng họ mô hình.
    \item Tính ma trận tương đồng cosine.
    \item Trích xuất Top-$k$ láng giềng gần nhất cho mỗi phim.
    \item Xuất artifact: \texttt{movie\_index.json}, \texttt{similar\_by\_movie.json}, \texttt{metadata.json}.
\end{enumerate}

\section{Mô hình nền}
\subsection{Baseline A: TF-IDF trên overview}
Mục tiêu của baseline A là tạo đường cơ sở tối giản với chi phí thấp. Mỗi phim được biểu diễn bằng vector TF-IDF trên trường \textit{overview}. Phương pháp này hiệu quả khi mô tả phim đủ dài và chứa từ khóa phân biệt thể loại.

\subsection{Baseline B: TF-IDF trên movie profile}
Baseline B giữ nguyên thuật toán TF-IDF nhưng mở rộng đầu vào sang profile tổng hợp. Đây là thí nghiệm quan trọng để lượng hóa mức đóng góp của review và genre metadata khi chưa thay đổi họ mô hình.

\section{Mô hình cải tiến}
\subsection{Sentence Transformer cho retrieval ngữ nghĩa}
Mô hình cải tiến sử dụng sentence embedding cho profile phim. So với TF-IDF, embedding dày đặc có khả năng gom cụm các cách diễn đạt khác nhau nhưng cùng nghĩa, nhờ đó thường cải thiện chất lượng truy hồi ở các truy vấn ``mô tả cảm xúc''.

\subsection{Chiến lược triển khai}
Thay vì fine-tune ngay từ đầu, đề tài đi theo lộ trình tăng dần:
\begin{enumerate}[label=\alph*.]
    \item inference với mô hình pretrained để thiết lập mốc chất lượng;
    \item phân tích lỗi theo nhóm truy vấn;
    \item chỉ fine-tune khi lợi ích kỳ vọng vượt chi phí huấn luyện.
\end{enumerate}

\section{Tách tập đánh giá}
Vì bài toán là retrieval theo item, cách chia tập truyền thống train/test trên mẫu độc lập không phản ánh đúng kịch bản triển khai. Đề tài áp dụng chia tập query--candidate theo hàm băm định danh phim (deterministic hash split), bảo đảm:
\begin{itemize}
    \item khả năng tái lập giữa các lần chạy;
    \item giảm phụ thuộc vào seed ngẫu nhiên;
    \item ngăn rò rỉ trực tiếp giữa truy vấn và tập ứng viên trong bước đánh giá.
\end{itemize}

\section{Độ phức tạp tính toán và chiến lược tối ưu}
Với $N$ phim, phép tính tương đồng đầy đủ có độ phức tạp $O(N^2)$. Đối với quy mô vài nghìn phim, cách làm này vẫn khả thi trong môi trường đồ án. Tuy nhiên, khi mở rộng quy mô, cần chuyển sang ANN indexing (Qdrant/Faiss) để giảm chi phí truy hồi online.

\section{Artifact-first serving}
Một quyết định kỹ thuật quan trọng của đề tài là phục vụ API từ artifact đã sinh offline thay vì huấn luyện trong runtime. Điều này giúp:
\begin{itemize}
    \item giảm rủi ro drift giữa môi trường huấn luyện và phục vụ;
    \item kiểm soát phiên bản mô hình rõ ràng;
    \item triển khai fallback về discovery nếu artifact chưa sẵn sàng.
\end{itemize}

\section{Mô tả thuật toán gợi ý}
\begin{lstlisting}[language=python,caption={Giả mã quy trình gợi ý phim tương đồng}]
for movie in movies:
    profile[movie] = build_profile(movie)

embeddings = encoder(profile)
sim_matrix = cosine_similarity(embeddings, embeddings)

for movie in movies:
    neighbors = top_k(sim_matrix[movie], k)
    save_recommendations(movie, neighbors)
\end{lstlisting}

\section{Kết luận chương}
Chương này trình bày toàn bộ lựa chọn phương pháp dưới dạng một quy trình nghiên cứu khả thi và có thể vận hành thật. Điểm cốt lõi là giữ nguyên đơn vị biểu diễn để so sánh công bằng các họ mô hình, đồng thời tách rõ huấn luyện và suy luận để đảm bảo tính ổn định khi tích hợp vào ứng dụng web.
